\documentclass[12pt]{article}
\usepackage[utf8]{inputenc}
\usepackage[T1]{fontenc}
\usepackage{hyperref}
\title{Adangbe Language Resources for K-12 Curriculum Development}
\author{Compiled by Claude Sonnet 5.5 (claude-sonnet-5-5) for Andr\'as Kornai}
\date{}
\begin{document}
\maketitle

\section{Introduction}
Adangbe (ISO 639-3 adq) is an endangered language of Ghana and Togo. Web search found very little material that is specific to it. Several results blur Adangbe with the much larger Adangme/Dangme language (ISO 639-3 ada), or describe it under the name Agotime. Check the ISO code before using any resource.

\section{Primary Educational Resources}
No primers, textbooks or school programmes for Adangbe turned up in search results.

\section{Linguistic and Reference Materials}
\begin{itemize}
\item Wikipedia, ``Agotime language'' (an encyclopedic overview; the sources do not fully agree on how Agotime and Adangbe are related): \url{https://en.wikipedia.org/wiki/Agotime_language}
\item Joshua Project language page for adq: \url{https://joshuaproject.net/index.php/languages/adq}
\end{itemize}
No dictionary or grammar specific to Adangbe was found.

\section{Religious and Community Resources}
\begin{itemize}
\item Scripture Earth page for adq (Bible portions and links): \url{https://www.scriptureearth.org/00eng.php?iso=adq}
\item Joshua Project, Adan/Adangbe in Togo: \url{https://joshuaproject.net/index.php/people_groups/10154/TO}
\end{itemize}

\section{Digital and Audio Resources}
\begin{itemize}
\item Global Recordings Network, language page for adq: \url{https://globalrecordings.net/en/language/adq}
\end{itemize}

\section{Curriculum Development Context}
Little documentation exists, and the language is not known to be taught in schools. A curriculum would first need fieldwork with the community: orthography, a lexicon and texts. Scripture portions and recordings are the only material found.
\end{document}